\chapter{Game Theory II: Dynamic Games and Commitment}
\label{chap:game-theory-dynamic}

\section{Learning objectives}

After completing this chapter, you should be able to represent sequential choices with game trees, use backward induction, define subgame-perfect equilibrium, and assess credible commitments.

\section{Sequential games}

In \keyterm{dynamic games}, players move in sequence and may observe earlier actions. A game tree records the order of moves, choices, and payoffs. \keyterm{Backward induction} solves a finite game by first determining what a player will do at the final decision node, then working back to earlier decisions.

\section{Credibility and subgame perfection}

A threat is \keyterm{credible} only if it is in the threatening player's interest to carry it out when the relevant moment arrives. A \keyterm{subgame-perfect equilibrium} is a strategy profile that specifies a Nash equilibrium in every subgame. It rules out non-credible threats.

Repeated interaction can support cooperation when future consequences are sufficiently important. But cooperation also depends on observability, patience, enforcement, and the possibility of punishment or forgiveness.

\section*{Chapter Review}

\subsection*{Key concepts}
\begin{itemize}
  \item \textbf{Game tree}: representation of sequential moves.
  \item \textbf{Backward induction}: solving from final choices backward.
  \item \textbf{Credible threat}: action optimal when the contingency occurs.
  \item \textbf{Subgame-perfect equilibrium}: equilibrium in every subgame.
\end{itemize}

\section*{Exercises}

\subsection*{Conceptual questions}
\begin{enumerate}[label=\arabic*.]
  \item Why can a Nash equilibrium in a dynamic game contain an incredible threat?
  \item Explain backward induction.
  \item How can repeated interaction support cooperation?
\end{enumerate}

\subsection*{Graphical exercises}
\begin{enumerate}[resume,label=\arabic*.]
  \item Draw a two-stage entry-deterrence game and solve it by backward induction.
\end{enumerate}

\subsection*{Analytical exercises}
\begin{enumerate}[resume,label=\arabic*.]
  \item An entrant chooses Enter or Stay Out. If Enter, an incumbent chooses Fight or Accommodate. Construct payoffs where Fight is not credible and find the subgame-perfect equilibrium.
  \item Explain why a costly commitment can change the outcome of a bargaining game.
\end{enumerate}

\subsection*{Exam-style multiple-choice questions}
\begin{enumerate}[resume,label=\arabic*.]
  \item Backward induction begins by analysing: (a) the first move; (b) the final decision nodes; (c) average payoffs; or (d) market demand.
  \item A subgame-perfect equilibrium excludes: (a) all cooperation; (b) credible actions; (c) non-credible threats; or (d) sequential decisions.
\end{enumerate}

\subsection*{PPE discussion question}

\noindent Political leaders and firms often make public commitments that limit their future freedom of action. When can such commitments improve cooperation or bargaining outcomes, and when can they be dangerous?
